\section{Poisson witnesses and their lower tail}
\label{sec:poisson}

The blocks constructed in Proposition~\ref{prop:blocks} give measurable
sets of slab normals, one set for each block of residuals. Choosing a
distinct Poisson point from each set certifies that all the residuals are
cut. We now estimate the probability that none of a finite collection of
targets has such a certificate. The lower-tail inequality originates in
the work of Janson~\cite{Janson1990LargeDeviations}. The continuous version
and diagram method below follow Li--Liu
\cite[Lemma~6.4, Equation~(6.16), and Lemma~9.6]{LiLiu2026Lattice}.
We include the finite
inequality, its passage to measurable slots, and the matching estimates
to specify all multiplicities and diagonal terms.

\subsection{The finite inequality}

The events in the finite approximation are indexed by ordered choices of
cells. Different choices can use the same set of cells, so repeated
underlying subsets must remain separate indices.

\begin{lemma}[Finite indexed Janson inequality]
\label{lem:finite-janson}
Let $(X_v)_{v\in V}$ be independent Bernoulli variables on a finite set
$V$, and let $(H_i)_{i=1}^N$ be an indexed family of nonempty subsets of
$V$, with repetitions allowed. Set
\[
 I_i=\prod_{v\in H_i}X_v,\qquad Y=\sum_{i=1}^NI_i,\qquad m=\mathbb{E}Y,
 \qquad
 \Delta^*=\sum_{\substack{1\le i,j\le N\\H_i\cap H_j\ne\varnothing}}
                  \mathbb{E}(I_iI_j).
\]
If $m>0$, then $\Delta^*\ge m$ and
\begin{equation}
 \mathbb{P}(Y=0)\le \exp\left(-\frac{m^2}{2\Delta^*}\right).
 \label{eq:finite-janson}
\end{equation}
The dependency sum is ordered and includes the diagonal.
\end{lemma}

\begin{proof}
We first recall Harris's correlation inequality for product Bernoulli
measures~\cite{Harris1960Correlation}, together with its elementary proof.
Increasing functions of independent Bernoulli variables have nonnegative
covariance. Indeed, induction on the number of variables and conditioning
on the last variable reduce the additional covariance term to
$p(1-p)(f_1-f_0)(g_1-g_0)\ge0$, where $f_b,g_b$ are the corresponding
conditional means. Thus an increasing event $B$ and a decreasing event
$F$ satisfy
\begin{equation}
 \mathbb{P}(B\mid F)\le\mathbb{P}(B)
 \quad\text{whenever }\mathbb{P}(F)>0.
 \label{eq:harris-opposite}
\end{equation}

Write $B_i=\{I_i=1\}$, $p_i=\mathbb{P}(B_i)$, and
$p_{ij}=\mathbb{P}(B_i\cap B_j)$. For each $i$, let $F_i$ be the event
that none of the earlier $B_j$ with $H_i\cap H_j=\varnothing$ occurs,
and let $G_i$ exclude the remaining earlier events. If some earlier
avoidance event has probability zero, the conclusion is immediate.
Otherwise all the conditional probabilities below are defined.
Since $F_i$ uses only coordinates outside $H_i$, it is independent of
$B_i$. It is also decreasing. A conditional union bound and
\eqref{eq:harris-opposite} therefore give
\[
 \begin{aligned}
 \mathbb{P}(B_i\mid F_i\cap G_i)
 &\ge\mathbb{P}(B_i\cap G_i\mid F_i)\\
 &\ge p_i-
       \sum_{\substack{j<i\\H_i\cap H_j\ne\varnothing}}p_{ij}.
 \end{aligned}
\]
The first inequality uses $\mathbb{P}(G_i\mid F_i)\le1$.
Multiplying the successive avoidance probabilities and using
$1-x\le e^{-x}$ yields
\begin{equation}
 \mathbb{P}(Y=0)\le\exp(-m+\Delta_{\mathrm{off}}/2),
 \qquad
 \Delta_{\mathrm{off}}=
 \sum_{\substack{i\ne j\\H_i\cap H_j\ne\varnothing}}p_{ij}.
 \label{eq:janson-first}
\end{equation}
This remains valid when an intermediate lower bound above is negative.

Attach to each index $i$ a private independent Bernoulli variable of
parameter $q$, and require it as part of $B_i$. The new mean is $qm$,
and its ordered off-diagonal sum is $q^2\Delta_{\mathrm{off}}$.
The original zero event implies the new zero event, so
\eqref{eq:janson-first} bounds its probability by
$\exp(-qm+q^2\Delta_{\mathrm{off}}/2)$. Since
$\Delta^*=m+\Delta_{\mathrm{off}}$, choosing $q=m/\Delta^*\in(0,1]$
proves \eqref{eq:finite-janson}. The proof retains every index even
when two supports coincide.
\end{proof}

Janson--Warnke \cite[Equation~(2)]{JansonWarnke2016LowerTail} record a
stronger constant in this indexed formulation. The constant in
Lemma~\ref{lem:finite-janson} suffices here.

\subsection{Measurable slots and matching diagrams}

Let $(\Omega,\mu)$ be a finite nonatomic standard Borel measure space,
and let $\Pi$ be a Poisson process of intensity $\mu$. Consider a finite
nonempty family of witness types indexed by $\alpha\in\mathcal A$.
In the covering application, a type is a target point and its slots are
blocks of residuals. For each $\alpha$, let $I_\alpha$ be a finite nonempty set of labeled
slots, with measurable sets $E_{\alpha i}\subseteq\Omega$ satisfying
$\lambda_{\alpha i}:=\mu(E_{\alpha i})>0$. Define its \emph{activity} by
\[
 \nu_\alpha=\prod_{i\in I_\alpha}\lambda_{\alpha i},
\]
and let $Z_\alpha$ count injective assignments of Poisson points to the
slots, with the point assigned to $i$ lying in $E_{\alpha i}$. These
are ordered counts: equal slot sets do not identify their labels.

A \emph{partial matching} between $I_\alpha$ and $I_\beta$ is a set
of pairs $(i,j)$ with no repeated first or second coordinate. Put
\begin{equation}
 D_{\alpha\beta}=
 \nu_\alpha\nu_\beta
 \sum_{\substack{M\text{ a partial matching}\\M\ne\varnothing}}
 \prod_{(i,j)\in M}
 \frac{\mu(E_{\alpha i}\cap E_{\beta j})}
      {\lambda_{\alpha i}\lambda_{\beta j}}.
 \label{eq:diagram-definition}
\end{equation}
The Poisson factorial-moment formula gives
\[
 \mathbb E Z_\alpha=\nu_\alpha,\qquad
 \mathbb E[Z_\alpha Z_\beta]
       =\nu_\alpha\nu_\beta+D_{\alpha\beta}.
\]
For the second identity, group pairs of injective assignments by the
partial matching of slots occupied by the same point. The empty matching
contributes $\nu_\alpha\nu_\beta$; the nonempty matchings give exactly
the terms in \eqref{eq:diagram-definition}. Thus $D_{\alpha\beta}$ is
the covariance of $Z_\alpha$ and $Z_\beta$.

\begin{proposition}[Poisson lower tail]
\label{prop:poisson}
Let $\Pi$ be a Poisson process on a finite nonatomic standard Borel
measure space. For each member of a finite nonempty family, take a
finite nonempty labeled list of measurable slots of positive intensity.
Let $Z_\alpha$ count their injective assignments, let $\nu_\alpha$ be
the product of their intensities, and define $D_{\alpha\beta}$ by
\eqref{eq:diagram-definition}. Then $\mathbb{E}Z_\alpha=\nu_\alpha$,
and
\begin{equation}
 \mathbb{P}(Z_\alpha=0\text{ for all }\alpha\in\mathcal A)
 \le\exp\left[-\frac{(\sum_\alpha\nu_\alpha)^2}
                        {2\sum_{\alpha,\beta}D_{\alpha\beta}}\right].
 \label{eq:poisson-janson}
\end{equation}
The denominator is positive: the full identity matching in
$D_{\alpha\alpha}$ contributes $\nu_\alpha$.
\end{proposition}

\begin{proof}
Write $\tau=\mu(\Omega)>0$. Conditional on $N=\Pi(\Omega)$, the
points are independent with law $\mu/\tau$ and are almost surely
distinct. If $k=|I_\alpha|$, the Poisson factorial moment
$\mathbb{E}(N)_k=\tau^k$ gives
\[
 \mathbb{E}Z_\alpha
 =\mathbb{E}(N)_k\,\tau^{-k}\prod_i\lambda_{\alpha i}
 =\nu_\alpha.
\]
In particular, no factorial divides the activity, even for repeated
slot sets.

We apply Lemma~\ref{lem:finite-janson} to cell-occupation indicators.
Fix $\delta>0$. Form the Boolean atoms generated by all the slot sets
and subdivide each positive-measure atom into finitely many measurable
cells of measure at most $\delta$, using nonatomicity. After omitting
null cells, this gives a finite partition $\mathcal P_\delta$ that
respects every slot. Modifications on the common null set do not affect
the Poisson process. For $C\in\mathcal P_\delta$, set
\[
 a_C=\mu(C),\qquad X_C=\mathbf1_{\{\Pi(C)>0\}},\qquad
 p_C=\mathbb{E}X_C=1-e^{-a_C}.
\]
These variables are independent and
\begin{equation}
 0\le a_C-p_C\le\tfrac12a_C^2\le\tfrac\delta2a_C.
 \label{eq:cell-occupation-error}
\end{equation}

For each $\alpha$, retain every injective assignment of cells
$f:I_\alpha\to\mathcal P_\delta$ with $f(i)\subseteq E_{\alpha i}$.
Its indicator is $\prod_iX_{f(i)}$. Let $Y_\delta$ be the sum over
these indexed assignments. Different assignments remain different
indices even when their underlying cell sets agree.

The following estimate controls all cell collisions. For a fixed list
$F_1,\ldots,F_r$ of unions of cells, let $Q_\delta(F_1,\ldots,F_r)$
be the sum of $\prod_{v=1}^rp_{C_v}$ over distinct cells
$C_v\subseteq F_v$. Then
\begin{equation}
 \begin{split}
 0\le\prod_v\mu(F_v)-Q_\delta(F_1,\ldots,F_r)
 \le\delta\bigg[&\frac r2\prod_v\mu(F_v)\\
 &+\sum_{v<w}\mu(F_v\cap F_w)
                      \prod_{u\ne v,w}\mu(F_u)\bigg].
 \end{split}
 \label{eq:cell-collision-error}
\end{equation}
Indeed, the unconstrained sum factors over the slots, and
\eqref{eq:cell-occupation-error}, followed by telescoping the product,
bounds its difference from $\prod_v\mu(F_v)$ by the first term.
The weight of assignments with $C_v=C_w$ is at most
\[
 \sum_{C\subseteq F_v\cap F_w}p_C^2
       \prod_{u\ne v,w}\mu(F_u)
 \le\delta\mu(F_v\cap F_w)\prod_{u\ne v,w}\mu(F_u).
\]
A union bound proves \eqref{eq:cell-collision-error}. Empty products
are 1, and the estimate also holds if a set has measure zero.
Applying it to each slot list gives
\begin{equation}
 m_\delta:=\mathbb{E}Y_\delta\longrightarrow\sum_\alpha\nu_\alpha.
 \label{eq:cell-mean-limit}
\end{equation}

It remains to identify the finite dependency sum and the limiting event.
For two internally injective assignments $f$ and $g$, the complete set
of equalities $f(i)=g(j)$ is a unique partial matching $M$.
Their Bernoulli supports intersect precisely when $M\ne\varnothing$.
For fixed $\alpha,\beta,M$, identify the slots paired by $M$ and
leave the other slots as singletons. Assign to each resulting class its
slot set, or the intersection of its two slot sets. Denote these sets
by $(F_v)_v$. Requiring distinct cells for these classes excludes
exactly the unintended equalities. Thus the contribution with exact
matching $M$ is $Q_\delta((F_v)_v)$, and its limit is
\[
 \prod_v\mu(F_v)=\nu_\alpha\nu_\beta
 \prod_{(i,j)\in M}
 \frac{\mu(E_{\alpha i}\cap E_{\beta j})}
      {\lambda_{\alpha i}\lambda_{\beta j}}.
\]
All sums are finite. Estimate~\eqref{eq:cell-collision-error} therefore
shows that the ordered dependency sum, including identical assignment
pairs, satisfies
\begin{equation}
 \Delta^*_\delta\longrightarrow\sum_{\alpha,\beta}D_{\alpha\beta}.
 \label{eq:cell-diagram-limit}
\end{equation}
Identical assignments give the full identity matching; other full
matchings caused by repeated slots remain separate terms.

Finally, couple the cell variables with the original process. The
probability that any cell contains two or more points is at most
\[
 \frac12\sum_C\mathbb{E}[\Pi(C)(\Pi(C)-1)]
 =\frac12\sum_Ca_C^2\le\frac{\delta\tau}{2}.
\]
Outside this event, point injections and occupied-cell injections are
in bijection, simultaneously for every witness type. Hence
\[
 \left|\mathbb{P}(Y_\delta=0)-
       \mathbb{P}(Z_\alpha=0\text{ for all }\alpha)\right|
 \le\frac{\delta\tau}{2}.
\]
Apply Lemma~\ref{lem:finite-janson} and let $\delta\downarrow0$,
using \eqref{eq:cell-mean-limit} and \eqref{eq:cell-diagram-limit}.
The identity-matching contribution makes the limiting denominator
positive. This proves \eqref{eq:poisson-janson}.
\end{proof}

The approximation is performed for each fixed finite family. Its size
and the total intensity may grow with dimension; no uniform estimate
in those quantities is needed before taking the cell limit. An empty
slot list has one empty injection and is treated as a deterministic
certificate before applying Proposition~\ref{prop:poisson}.

\subsection{Self and cross diagrams}

The two budgets in Proposition~\ref{prop:blocks} now control the
self-diagram sum. The factor $1+\lambda_i^{-1}$ pays for identifying
the same slot in both copies of a witness.

\begin{lemma}[Matching bounds]
\label{lem:diagrams}
For positive-measure slot lists and the diagrams in
\eqref{eq:diagram-definition}, the following bounds hold.
For one witness type, suppress its label and set
\[
 \delta_i=1+\lambda_i^{-1},\qquad
 w_{ij}=\frac{\mu(E_i\cap E_j)}
 {\sqrt{\lambda_i(1+\lambda_i)\lambda_j(1+\lambda_j)}}
 \quad(i\ne j).
\]
Then
\begin{equation}
 1+\frac{D_{\alpha\alpha}}{\nu_\alpha^2}
 \le\prod_i\delta_i\prod_i\left(1+\sum_{j\ne i}w_{ij}\right).
 \label{eq:self-diagram-product}
\end{equation}
For two witness types, put
\[
 \rho_{\alpha\beta}=
 \sum_{i\in I_\alpha,j\in I_\beta}
 \frac{\mu(E_{\alpha i}\cap E_{\beta j})}
      {\lambda_{\alpha i}\lambda_{\beta j}}.
\]
Then
\begin{equation}
 D_{\alpha\beta}\le\nu_\alpha\nu_\beta
                         (e^{\rho_{\alpha\beta}}-1).
 \label{eq:cross-diagram-product}
\end{equation}
In particular, if every cross-slot intersection has measure zero,
then $D_{\alpha\beta}=0$.
\end{lemma}

\begin{proof}
Write $h_{ij}=\mu(E_i\cap E_j)/(\lambda_i\lambda_j)$.
For the self sum, fix the off-diagonal edges $P$ of a partial matching
between two copies of the slot list. Regard $P$ as a directed graph on
the common label set. Every vertex has total degree
$d_P(i)\in\{0,1,2\}$, since its in-degree and out-degree are at most
one. A diagonal loop is available exactly at the vertices of degree
zero. Since $h_{ii}=\lambda_i^{-1}$, summing the compatible loop
choices gives
\[
 \left(\prod_i\delta_i\right)
 \left(\prod_{(i,j)\in P}h_{ij}\right)
 \left(\prod_{i:d_P(i)>0}\delta_i^{-1}\right)
 \le\left(\prod_i\delta_i\right)\prod_{(i,j)\in P}w_{ij}.
\]
Here $\delta_i^{-1}\le \delta_i^{-d_P(i)/2}$ at each incident vertex.
This includes directed paths and cycles; Figure~\ref{fig:self-matching}
illustrates the permitted incidences. Summing over $P$ and dropping
the requirement of distinct right endpoints bounds the remaining sum
by $\prod_i(1+\sum_{j\ne i}w_{ij})$. The empty matching contributes
1, proving \eqref{eq:self-diagram-product}.

For two lists, drop all matching restrictions and sum over nonempty
subsets of cross edges. With the analogous nonnegative weights $h_{ij}$,
the result is at most
$\prod_{i,j}(1+h_{ij})-1\le e^{\sum_{i,j}h_{ij}}-1$.
This proves \eqref{eq:cross-diagram-product}.
\end{proof}

\begin{figure}[htbp]
\centering
\begin{tikzpicture}[
  x=1cm,y=1cm,
  slot/.style={circle,draw=black!70,line width=.6pt,
    minimum size=8mm,inner sep=2pt,font=\small},
  edge/.style={-{Stealth[length=2.2mm,width=1.5mm]},
    draw=blue!45!black,line width=.8pt},
  note/.style={font=\small,align=center,text=black!80}
]
  \node[slot] (one) at (0,0) {$1$};
  \node[slot] (two) at (2,0) {$2$};
  \node[slot] (three) at (4,0) {$3$};
  \draw[edge] (one) -- (two);
  \draw[edge] (two) -- (three);
  \node[note] at (2,1.35) {directed path};
  \node[note] at (0,-.8) {$d_P(1)=1$};
  \node[note] at (2,-.8) {$d_P(2)=2$};
  \node[note] at (4,-.8) {$d_P(3)=1$};

  \node[slot] (four) at (6.5,0) {$4$};
  \node[slot] (five) at (8.5,0) {$5$};
  \draw[edge] (four) to[bend left=27] (five);
  \draw[edge] (five) to[bend left=27] (four);
  \node[note] at (7.5,1.35) {directed $2$-cycle};
  \node[note] at (6.5,-.8) {$d_P(4)=2$};
  \node[note] at (8.5,-.8) {$d_P(5)=2$};

  \node[slot] (six) at (11.5,0) {$6$};
  \draw[edge,dashed] (six) to[out=135,in=45,looseness=5] (six);
  \node[note] at (11.5,1.35) {optional diagonal match};
  \node[note] at (11.5,-.8) {$d_P(6)=0$};
\end{tikzpicture}
\caption{A self-matching after identifying the two copies of the slot labels.
The solid edges form the off-diagonal matching
$P=\{(1,2),(2,3),(4,5),(5,4)\}$; each label has at most one incoming
and one outgoing edge. The degree $d_P(i)$ counts both.
A diagonal match is available only at an isolated label, as illustrated
by the optional dashed loop at $6$; even the degree-one path endpoints
cannot receive a diagonal match.}
\label{fig:self-matching}
\end{figure}


\subsection{Application to residual caps}

We return to the fixed parameters and the spherical intensity $\mu$
of Section~\ref{sec:construction}. The following consequence collects
the inputs needed for the target selection in the next section.

\begin{corollary}[Witnesses for residual lists]
\label{cor:witnesses}
Fix $A_0>0$. There are constants $A,C_*,H>0$ such that, for all
sufficiently large $n$, the following holds. Let $R\ge1$, and let
$(x_{\alpha i})_{i=1}^{m_\alpha}$, $\alpha\in\mathcal A$, be any
finite nonempty family of deterministic lists satisfying
\[
 1\le m_\alpha\le n^2,\qquad
 r_--3\eta\le |x_{\alpha i}|\le b+3\eta.
\]
Define
\[
 C_{\alpha i}=\{u\in S^{n-1}:|u\cdot x_{\alpha i}|>1+3\eta\},
 \qquad
 s_{\alpha i}=-\log\min\{\mu(C_{\alpha i}),1/2\},
\]
and suppose $m_\alpha+\sum_i s_{\alpha i}\le A_0R$.
For each list there is a nonempty block partition and measurable slots
$E_{\alpha S}\subseteq\bigcap_{i\in S}C_{\alpha i}$, chosen
independently of $\Pi$, with positive intensities at most $1/2$.
Their activities and diagrams satisfy
\begin{equation}
 e^{-AR}\le\nu_\alpha\le1,\qquad
 D_{\alpha\alpha}\le\nu_\alpha^2e^{C_*R}.
 \label{eq:witness-self}
\end{equation}
For distinct $\alpha,\beta$, if all cross-residual angles exceed
$H\sqrt{\log n/n}$, then
\begin{equation}
 D_{\alpha\beta}\le n^{-3}\nu_\alpha\nu_\beta.
 \label{eq:witness-cross}
\end{equation}
If all those angles exceed $4\sqrt\varepsilon$, then
$D_{\alpha\beta}=0$. A positive injection count for a list implies
that every cap in that list is hit by $\Pi$.
\end{corollary}

\begin{proof}
Apply Lemma~\ref{lem:caps} with shell parameter $10$ and perturbation
bound $3$. Its affinity and common-intersection estimates supply the
hypotheses of Proposition~\ref{prop:blocks}; the common-intersection
constant is available before the radius multiplier in that proposition
is fixed. Apply the proposition separately to each list, making all
choices before sampling $\Pi$. The budget
\eqref{eq:block-first-budget} gives
\[
 -\log\nu_\alpha
 \le\sum_S\log(1+\lambda_{\alpha S}^{-1})\le AR.
\]
Together with \eqref{eq:block-row-budget} and
\eqref{eq:self-diagram-product}, it also gives the self bound in
\eqref{eq:witness-self}, after choosing $C_*$.

To check the cross estimate, write
$\bar c_{\alpha i}=\min\{\mu(C_{\alpha i}),1/2\}$ and let
$i(S)$ be the anchor of a block. Lemma~\ref{lem:caps} and
Proposition~\ref{prop:blocks} provide fixed exponents $P_0,P_1>0$
such that
\[
 \bar c_{\alpha i}\ge n^{-P_0},\qquad
 \lambda_{\alpha S}\ge n^{-P_1}\bar c_{\alpha,i(S)},\qquad
 E_{\alpha S}\subseteq C_{\alpha,i(S)}.
\]
Apply the affinity estimate to each two-element list of anchor residuals;
this does not require the union of the two full lists to have at most
$n^2$ entries. It gives
\[
 \frac{\mu(E_{\alpha S}\cap E_{\beta T})}
      {\lambda_{\alpha S}\lambda_{\beta T}}
 \le Cn^{2P_1+P_0}
       e^{-c n\angle(x_{\alpha,i(S)},x_{\beta,i(T)})^2}.
\]
There are at most $n^4$ block pairs. Choose $H$ after these constants,
so that $cH^2>9+2P_1+P_0$. For sufficiently large $n$, the angular
hypothesis gives $\rho_{\alpha\beta}\le n^{-4}$. Equation
\eqref{eq:cross-diagram-product} now gives
$D_{\alpha\beta}\le2n^{-4}\nu_\alpha\nu_\beta$, which implies
\eqref{eq:witness-cross}. At cross angles greater than
$4\sqrt\varepsilon$, Corollary~\ref{cor:caps-disjoint}
makes every cross-slot intersection empty, and the diagram is zero.

Finally, the point assigned to a block lies in every cap belonging to
that block. A positive injection count thus hits every residual cap.
\end{proof}

For a finite center list $(p_i)$ and a target $y$, take all residuals
$x_i=y-p_i$ with $|x_i|\le b+3\eta$. If this list satisfies the
hypotheses of Corollary~\ref{cor:witnesses}, a positive witness excludes
each such residual from $K^+$ by a strict slab inequality. All remaining
residuals are excluded by its outer ball. Consequently
\begin{equation}
 Z_y>0\quad\Longrightarrow\quad
 y\notin\bigcup_i(p_i+K^+).
 \label{eq:witness-hole}
\end{equation}
If the relevant list is empty, the target is already a hole.
Proposition~\ref{prop:poisson} supplies an exponential bound on the
probability that every chosen target lacks a witness. A second-moment
bound alone would not give the decay required for the finite-pattern
union in the final argument.
